\SourceLeaf{239}%
einsetzen und berücksichtigen, dass die sich so ergebenden Werte der Differentiale \(\mathrm{d}t,\mathrm{d}u\) unabhängig von den Differentialen \(\mathrm{d}x,\mathrm{d}y\) jeweils den Größen \(T\mathrm{d}x+U\mathrm{d}y\) und \(U\mathrm{d}x+V\mathrm{d}y\) gleich sein müssen, so werden wir nach einigen recht naheliegenden Umformungen Folgendes finden:
\[
\begin{aligned}
C^3T&=\alpha Ab'b'+\beta Bb'b'+\gamma Cb'b'\\
&\quad-2\alpha'Abb'-2\beta'Bbb'-2\gamma'Cbb'\\
&\quad+\alpha''Abb+\beta''Bbb+\gamma''Cbb\\
C^3U&=-\alpha Aa'b'-\beta Ba'b'-\gamma Ca'b'\\
&\quad+\alpha'A(ab'+ba')+\beta'B(ab'+ba')+\gamma'C(ab'+ba')\\
&\quad-\alpha''Aab-\beta''Bab-\gamma''Cab\\
C^3V&=\alpha Aa'a'+\beta Ba'a'+\gamma Ca'a'\\
&\quad-2\alpha'Aaa'-2\beta'Baa'-2\gamma'Caa'\\
&\quad+\alpha''Aaa+\beta''Baa+\gamma''Caa
\end{aligned}
\]
Setzen wir also der Kürze halber
\[
\begin{aligned}
A\alpha+B\beta+C\gamma&=D\qquad .\ .\ .\ .\ .\ .\ .\ .\ .\ .\quad(1)\\
A\alpha'+B\beta'+C\gamma'&=D'\qquad .\ .\ .\ .\ .\ .\ .\ .\ .\ .\quad(2)\\
A\alpha''+B\beta''+C\gamma''&=D''\qquad .\ .\ .\ .\ .\ .\ .\ .\ .\ .\quad(3)
\end{aligned}
\]
so ergibt sich
\[
\begin{aligned}
C^3T&=Db'b'-2D'bb'+D''bb\\
C^3U&=-Da'b'+D'(ab'+ba')-D''ab\\
C^3V&=Da'a'-2D'aa'+D''aa
\end{aligned}
\]
Durch Ausmultiplizieren finden wir daraus
\[
C^6(TV-UU)=(DD''-D'D')(ab'-ba')^2=(DD''-D'D')CC
\]
und damit die Formel für das Krümmungsmaß
\[
k=\frac{DD''-D'D'}{(AA+BB+CC)^2}
\]
\Article{11}
Auf der eben gefundenen Formel bauen wir nun eine weitere auf, die zu den fruchtbarsten Sätzen in der Theorie der gekrümmten Flächen zu zählen ist. Führen wir die folgenden Bezeichnungen ein:
